\section{问题提出的背景}

\subsection{背景介绍}

随着虚拟现实、增强现实、数字孪生、空间计算以及具身智能等方向持续发展，真实场景的数字化表达已不再满足于简单记录图像，而是希望进一步获得可渲染、可测量、可交互的三维场景表示。
与传统二维影像只能从固定机位回看不同，三维重建方法需要同时恢复场景的几何结构、外观信息和相机位姿，从而支持自由视点浏览、空间理解、交互仿真以及后续的任务规划与算法验证。

在这一背景下，以 NeRF \cite{mildenhall2021nerf} 和三维高斯泼溅（3DGS）\cite{kerbl2023} 为代表的新视角合成方法推动了三维场景表示的发展。
其中，3DGS 由于兼顾较高的重建质量与较快的训练、渲染效率，已经成为当前高保真场景重建的重要技术路线。
然而，现有 3DGS 方法大多建立在透视相机与较充分多视角采样的前提下，对于真实室内场景中普遍存在的遮挡频繁、弱纹理区域多、重复纹理明显以及采集轨迹受限等问题，仍然缺乏足够稳健的解决方案。

对于单个孤立物体的重建，研究者往往可以通过围绕目标拍摄 360 度多视角图像，获得较完整的角度覆盖，因此模型能够较容易地恢复稳定且细致的几何结构。
但在场景级重建任务中，这种理想采集条件通常难以满足。
受空间布局、遮挡关系、采集时间和设备运动轨迹限制，场景中的许多区域只能被从有限角度观测，未被充分覆盖的位置容易在重建结果中表现为模糊、拉伸、漂浮伪影或结构错位。
因此，如何在视角覆盖不足的条件下仍然获得具有较好几何一致性的高保真重建结果，已成为当前 3DGS 研究中的关键问题之一。

\subsubsection{项目提出的原因}

尽管 3DGS 在新视角合成质量和实时渲染效率方面表现突出，但其优化过程在本质上仍然主要依赖训练图像上的光度监督 \cite{kerbl2023}。当
输入视角数量不足、分布不均或局部区域长期处于弱观测状态时，场景恢复就会呈现出明显的欠定特征，即存在多组不同的三维结构都可能在已观测视角上获得相近的拟合误差。
此时，模型容易优先选择能够解释训练图像的解，而不一定是真实几何上最合理的解，进而在外插视角或遮挡边界处出现漂浮伪影、几何漂移、纹理拉伸和结构不连续等问题。

围绕这些难点，当前提高 3DGS 重建质量的研究已经从单纯提升像素指标，逐步转向引入更强的先验和约束。
一类工作通过抗锯齿与尺度一致性设计改善跨分辨率渲染稳定性，如 Mip-Splatting \cite{yu2024mipsplatting}；一类工作通过改变高斯表示或增加法线、深度等约束来强化几何准确性，如 2DGS \cite{huang2024twodgs}；还有一些方法尝试重新组织高斯的表示与优化方式，例如通过基于锚点的结构化表示和视角自适应高斯生成来提升表示效率与优化稳定性，Scaffold-GS 便是其中较具代表性的工作 \cite{lu2024scaffold}；
此外，也有方法通过跨视角深度一致性、透明度解耦或更明确的几何建模来抑制漂浮伪影和几何伪影，如 StableGS 与 Geometry-Grounded GS \cite{wang2025stablegs,zhang2026geometrygroundedgs}。
这些工作的共同目标，都是通过正则化、结构先验或模型约束来缩小解空间，提高在稀疏观测条件下的优化稳定性。

但从问题本质上看，上述方法多数仍然是在现有 RGB 观测基础上，通过更强的模型假设来约束优化过程。
这样的思路能够在一定程度上改善结果，却也可能引入较强的归纳偏置，使模型在复杂场景、特殊结构或异常观测条件下的表达能力受到限制。
换言之，如果输入数据本身对场景几何的覆盖仍然不充分，那么单纯依赖模型端正则化未必能够彻底解决欠定问题。
因此，本研究尝试从数据层面引入额外几何约束：一方面，鱼眼相机具有更大的视场角，能够在单张图像中覆盖更广的空间区域，为同一区域提供更多交叉观测关系；
另一方面，雷达可以提供与纹理、光照相对解耦的稀疏几何信息，为 3DGS 优化提供独立于 RGB 的结构锚点，从而提升跨视角的几何一致性。

然而，要真正发挥这两类信息源的优势，还存在若干关键技术障碍。
传统 3DGS 管线通常只面向透视相机建模，难以直接适配鱼眼成像下的强畸变投影；鱼眼相机模型参数较多，内参、畸变参数与相机位姿之间存在较强耦合，在运动恢复结构（SfM）标定和位姿恢复过程中若直接联合优化过多自由度，容易对局部匹配误差与观测噪声产生过拟合，从而导致估计得到的相机内参与外参不够准确，并进一步影响后续重建结果的几何一致性与整体质量；
与此同时，雷达点云本身存在稀疏、噪声大和离群点多等特点，如何对其进行有效筛选，并将其转换为适合 3DGS 优化的几何约束形式，同样是尚待解决的重要问题。基于以上原因，本课题提出融合鱼眼相机与雷达几何约束的室内场景重建方法，试图在稀疏观测条件下提升 3DGS 的几何稳定性与重建质量。

\subsection{本研究的意义和目的}

\subsubsection{研究目的}

本研究面向场景级 3DGS 重建中因视角覆盖不充分而导致几何退化这一核心问题，拟构建一个融合鱼眼图像与雷达的高保真室内场景重建框架。具体而言，本研究希望突破传统 3DGS 仅支持透视相机的限制，使其能够直接处理鱼眼相机成像模型下的图像观测；
在此基础上，进一步研究如何利用经过筛选的雷达点云为重建过程提供额外的几何锚点和一致性约束，从而降低模型仅依赖 RGB 光度误差进行拟合时出现的伪解风险。

从任务目标上看，本研究并不只是追求在已观测视角上的图像拟合能力，而是更关注在有限采样条件下场景几何结构的稳定恢复能力。通过引入鱼眼相机更大的视场角，
本研究希望在不显著增加采集负担的情况下提升单帧观测的信息密度，缓解场景局部区域观测不足的问题；
通过引入雷达提供的几何信息，本研究希望增强跨视角之间的结构一致性，抑制重建中常见的漂浮伪影、局部错位和几何变形。最终目标是建立一套在稀疏、多遮挡、弱纹理室内环境下仍具有较好鲁棒性的 3DGS 重建方法。

\subsubsection{研究意义}

本研究的意义首先体现在方法层面。当前针对 3DGS 的改进，大多集中在通过正则化、表面先验或模型结构调整来缩小可行解空间，本质上是在算法端为欠定问题引入附加假设。
这样的思路虽然有效，但约束越强，模型受到的归纳偏置也可能越明显。
相比之下，本研究尝试从数据层面补充约束，通过鱼眼相机扩大视野、通过雷达增强几何监督，使模型在保留较强表达能力的同时获得更可靠的优化依据。
这为如何在不显著牺牲模型自由度的前提下提升 3DGS 几何质量提供了一种新的研究思路。

其次，本研究具有明确的工程应用价值。
对于单物体重建而言，围绕目标进行高密度、多角度拍摄通常是可行的；
但在真实室内场景中，要对每个物体都获得完整的角度覆盖往往代价很高。
鱼眼相机能够以较少的图像数量覆盖更大范围空间，天然适合场景级快速采集；
雷达则能够在纹理缺失、光照变化较大或表面重复纹理较强的区域提供稳定的几何参考。
因此，若本研究能够实现两者与 3DGS 的有效融合，将有望降低室内场景高保真建模对密集采样轨迹的依赖，提高三维数字化系统在数字孪生、机器人感知、空间计算和仿真环境构建中的实用性。

最后，本研究也具有一定的前沿探索意义。
它不仅涉及鱼眼成像模型下的 3DGS 表达与优化问题，还涉及多模态传感器之间的协同建模、几何约束设计以及噪声点云的有效利用问题。
对这些问题的研究，有助于推动 3DGS 从以 RGB 为核心的单模态重建框架发展到具备多源几何信息融合能力的重建框架，并为后续进一步结合深度、激光或其他传感器信息提供可借鉴的方法基础。
